\begin{abstract}
Large-scale deep learning systems increasingly serve heterogeneous query streams by dynamically blending specialized parameter-efficient experts. While dynamic routing layers are central to these architectures, existing designs suffer from a fundamental, unresolved bottleneck: stateless routers adapt rapidly but exhibit extreme high-frequency routing jitter (noise) under observation perturbations, whereas stateful alternatives (e.g., biochemical ODEs or exponential filters) smooth the ensembling trajectory but introduce severe representational lag (inertial drag) at task boundaries. To resolve this dilemma, we propose \textbf{Active Inference Routing (AIR)}, a brain-inspired paradigm shift that models the multi-expert routing layer as an active, self-organizing cognitive agent performing test-time perception and action. Rather than relying on feedforward heuristics or static temporal filters, AIR represents the active task context as a stateful Gaussian variational belief, updating the ensembling weights through the minimization of Variational Free Energy. This formulation elegantly decomposes into precision-weighted sensory and prior prediction errors, enabling the router to dynamically balance top-down expectations with bottom-up stimuli. Under static variational covariance, this exact closed-form belief update is mathematically equivalent to a classical linear state observer (Kalman filter), providing a first-principles variational derivation of linear filtering for modular deep networks. Evaluated on the Analytical Coordinate Sandbox, AIR matches the theoretical representation alignment ceilings of an omniscient expert oracle under rapid, non-stationary task transitions. Meanwhile, under a stable but noisy stream paradigm, AIR simultaneously stabilizes ensembling trajectories, slashing routing jitter by up to 2.49$\times$. Finally, a mechanistic ablation study confirms that inhibitory pathways (negative feedback coupling) in the generative mapping are mathematically mandatory to actively suppress obsolete task representations and eliminate lag, offering a rich and theoretically grounded foundation for next-generation adaptive serving.
\end{abstract}
